\documentclass{article}
\usepackage{graphicx} % Required for inserting images

\title{LeNet-5 Paper Reproduction: Notes}
\author{Ilagaba Shehe}
\date{February 28, 2026}

\begin{document}

\maketitle

\section*{Section 2: Convolutional Neural Networks for
Isolated Character Recognition}
\begin{itemize}
    \item Convolutional Networks force the extraction of local features by restricting the receptive fields of hidden units to be local.
    \item Convolutional Networks combine three architectural ideas to ensure some degree of shift, scale, and distortion invariance, local receptive fields, shared weights (or weight replication) and spatial or temporal subsampling.
    \item Units in a layer are organized in planes within which all the units share the same set of weights. The set of outputs of the units in such a plane is called a feature map. Units in a feature map are all constrained to perform the same operation on different parts of the image. A complete convolutional layer is composed of several feature maps (with different weight vectors), so that multiple features can be extracted at each location.
    \item A simple way to reduce the precision with which the position of distinctive features are encoded in a feature map is to reduce the spatial resolution of the feature map. This can be achieved with a so called subsampling layers which performs a local averaging and a subsampling, reducing the resolution of the feature map, and reducing the sensitivity of the output to shifts and distortions.
    \item Fixed size convolutional networks that share weights along a single temporal dimension are known as Time Delay Neural Networks (TDNNs).
    
\end{itemize}
    \subsection*{LeNet-5 Architecture}
    \begin{itemize}
        \item The network comprises of 7 layers (not counting the input). They call contain trainable parameters (weights).
        \item the input is a 32x32 pixel images. Which are normalized
        \item Receptive Field(s): It is the size of the input area that a particular CNN feature is "looking at".
        \item Convolutional layers are labeled $C_x$ and sub-sampling layers are labeled $S_x$ and fully-connected layers are labeled $F_x$, where x is the layer index.
            \begin{itemize}
                \item Layer $C_1$: has 6 feature maps with size 28x28, where each feature map is connected to a 5x5 neighborhood in the input. $C_1$ contains 156 trainable parameters, and 122,304 connections.
                \item Layer $S_2$: has 6 feature maps of size 14x14. Each unit connected to a 2x2 neighborhood corresponding to the feature map in $C_1$. Uses a sigmoidal activation function. $S_2$ has 12 trainable parameters and 5,880 connections.
                \item Layer $C_3$: has 16 feature maps. Each unit in each feature map is connected to several 5x5 neighborhoods at identical locations in a subset of $S_2$'s feature maps. Layer $C_3$ has 1,516 trainable parameters and 151,600 connections.
                \item Layer $S_4$: has 16 feature maps of size 5x5. Each unit in each feature map is connected to a 2x2 neighborhood in the corresponding feature map in $C_3$ like in $C_1$ and $S_2$. $S_4$ has 32 trainable parameters and 2,000 connections.
                \item Layer $C_5$: has 120 feature maps where each unit is connected to a 5x5 neighborhood on all 16 f $S_4$'s feature maps. Layer $C_5$ has 48,120 trainable connections.
                \item Layer $F_6$: contains 84 units and is fully connected to $C_5$, and has 10,164 trainable parameters.
            \end{itemize}
        \item Like classical neural networks, units in layers up to $F_6$ compute a dot product between their input vector and their weight vector to which a bias is added. 
            \[y = w^Tx + b\]
        Then a that sum ($y$) is then input for a sigmoid function to produce the state for one of units in the next layer. For the CNN is squashing function is a scaled hyperbolic tangent:
            \[f(a) = Atanh(Sa)\]
        where A is the amplitude of the function and S determines its slope at the origin. A is chosen to be 1.7159.
        \item The output layer is composed f Euclidean Radial Basis Function units (RBF), one for each class, with 84 inputs each. The outputs of each RBF unit $y_i$ is computed as follows:
            \[y_i - \sum_j (x_j - w_{ij})^2 \]
        essentially each output RBF unit computes the Euclidean distance between its input vector and is parameter vector.
        \item Loss Function
            \begin{itemize}
                \item The LeNet-5 uses the Maximum Likelihood Estimation criterion (MLE).
                    \[E(W) = \frac{1}{P} \sum_{p=1}^{P} y_{D_p}(Z^P, W)\]
                Where $y_{D_p}$ is the output of the $D_p$-th RBF unit, which is the oen that corresponds to the correct class of input pattern $Z^p$. However there are three issues with this equation, where it was modified into this loss function equation:
                    \[E(W) = \frac{1}{P} \sum_{p=1}^{P} y_{D_p}(Z^P, W) + log(e^{-j} + \sum_i e^{-y_i(Z^p, W)})\].
                \item The log term represents the competitive nature of the classification. It forces the network to not only decrease the error for the correct class ( $_{D_P}$) but also increase the error (distance) for all other incorrect classes.
                \item Computing the gradient of the loss function with for all the weights of the layers in a convolutional network is done with back-propagation.
            \end{itemize}
    \end{itemize}

\end{document}
